% ===== PRÉAMBULE CANONIQUE — à coller VERBATIM en tête de CHAQUE .tex =====
\documentclass[11pt,a4paper]{article}
\usepackage[utf8]{inputenc}\usepackage[T1]{fontenc}\usepackage{lmodern}
\usepackage[french]{babel}
\usepackage{amsmath,amssymb,mathtools}\usepackage{icomma}
\usepackage{siunitx}\sisetup{locale=FR}  % \qty{}{}, \si{}, \ang{}
\usepackage{xcolor}\usepackage{colortbl}
\usepackage{tikz}\usepackage{pgfplots}\pgfplotsset{compat=1.18}
\usepackage[siunitx,europeanresistors]{circuitikz} % circuits (élec.)
\usepackage[most]{tcolorbox}\usepackage{enumitem}\usepackage{array,booktabs}
\usepackage[a4paper,margin=2.2cm]{geometry}
\usetikzlibrary{arrows.meta,calc,angles,quotes,positioning,patterns,%
  backgrounds,shapes.geometric,decorations.pathmorphing,decorations.markings,intersections}
\definecolor{primary}{RGB}{222,49,99}\definecolor{primarydark}{RGB}{150,22,55}
\definecolor{defcol}{RGB}{0,114,178}\definecolor{methcol}{RGB}{52,92,148}
\definecolor{excol}{RGB}{0,135,90}\definecolor{attncol}{RGB}{200,40,55}
\tcbset{boxbase/.style={enhanced,breakable,boxrule=0.9pt,arc=2.5pt,
  left=3mm,right=3mm,top=2.3mm,bottom=2.3mm,fonttitle=\bfseries,coltitle=white}}
% --- boîtes TUTORIEL (noms EXACTS, ne pas renommer) ---
\newtcolorbox{cle}[1][]{boxbase,colback=defcol!6,colframe=defcol,
  title={\textsc{À retenir}\ifx\relax#1\relax\relax\else\textmd{ : \itshape #1}\fi}}
\newtcolorbox{methode}[1][]{boxbase,colback=methcol!7,colframe=methcol,sharp corners,
  title={\textsc{Méthode}\ifx\relax#1\relax\relax\else\textmd{ --- \itshape #1}\fi}}
\newtcolorbox{exemplebox}[1][]{boxbase,colback=excol!5,colframe=excol,
  title={\textsc{Exemple}\ifx\relax#1\relax\relax\else\textmd{ : \itshape #1}\fi}}
\newtcolorbox{piege}[1][]{boxbase,colback=attncol!6,colframe=attncol,
  title={\textsc{Piège}\ifx\relax#1\relax\relax\else\textmd{ : \itshape #1}\fi}}
\newtcolorbox{remarque}[1][]{boxbase,colback=black!4,colframe=black!45,
  title={\textsc{Remarque}\ifx\relax#1\relax\relax\else\textmd{ : \itshape #1}\fi}}
% --- boîtes EXERCICE / CORRIGÉ (noms EXACTS, auto-comptées) ---
\newtcolorbox[auto counter]{exo}[1][]{boxbase,colback=primary!4,colframe=primary,
  title={\textsc{Exercice}~\thetcbcounter\ifx\relax#1\relax\relax\else\textmd{ --- \itshape #1}\fi}}
\newtcolorbox[auto counter]{corrige}[1][]{boxbase,colback=excol!4,colframe=excol,
  title={\textsc{Corrigé}~\thetcbcounter\ifx\relax#1\relax\relax\else\textmd{ --- \itshape #1}\fi}}
\newcommand{\vv}[1]{\vec{#1}}\newcommand{\ds}{\displaystyle}
\newcommand{\R}{\mathbb{R}}\newcommand{\N}{\mathbb{N}}\newcommand{\Z}{\mathbb{Z}}
% =========================================================================
\begin{document}

\begin{exo}[tout relier : le diamant]
Une lumière verte de longueur d'onde $\lambda_0=\qty{500}{nm}$ dans le vide arrive de l'air
($n_{\text{air}}=1$) sur un diamant ($n=2{,}42$) sous un angle d'incidence $\hat{\imath}_1=\ang{40}$.
On prend $c=\qty{3.0e8}{m/s}$.
\begin{enumerate}[label=\alph*)]
  \item Calcule la célérité de la lumière dans le diamant.
  \item Calcule la longueur d'onde de cette lumière dans le diamant.
  \item Calcule l'angle de réfraction $\hat{\imath}_2$ dans le diamant.
  \item Vérifie que $\dfrac{\sin\hat{\imath}_1}{\sin\hat{\imath}_2}$, $\dfrac{n}{n_{\text{air}}}$,
        $\dfrac{c_{\text{air}}}{c_{\text{diamant}}}$ et $\dfrac{\lambda_{\text{air}}}{\lambda_{\text{diamant}}}$
        donnent bien la même valeur.
\end{enumerate}
\end{exo}

\begin{exo}[identifier un milieu par sa longueur d'onde]
Une lumière monochromatique a une longueur d'onde $\lambda_1=\qty{650}{nm}$ dans un milieu 1 et
$\lambda_2=\qty{500}{nm}$ dans un milieu 2.
\begin{enumerate}[label=\alph*)]
  \item Calcule le rapport des célérités $c_1/c_2$.
  \item Le milieu 1 est de l'air ($n_1=1$). Détermine l'indice $n_2$ du milieu 2, et dis lequel des
        deux milieux est le plus réfringent.
\end{enumerate}
\end{exo}

\begin{exo}[remonter la chaîne depuis les angles]
À la traversée d'un dioptre, on mesure un angle d'incidence $\hat{\imath}_1=\ang{45}$ (milieu 1) et
un angle de réfraction $\hat{\imath}_2=\ang{28}$ (milieu 2).
\begin{enumerate}[label=\alph*)]
  \item Calcule le rapport $\dfrac{n_2}{n_1}$. Que valent alors $\dfrac{c_1}{c_2}$ et
        $\dfrac{\lambda_1}{\lambda_2}$ ?
  \item Dans quel milieu la lumière va-t-elle le plus vite ? Lequel est le plus réfringent ?
  \item Si la longueur d'onde vaut $\lambda_1=\qty{600}{nm}$ dans le milieu 1, que vaut $\lambda_2$ ?
\end{enumerate}
\end{exo}

\end{document}
